% Part: normal-modal-logic
% Chapter: completeness
% Section: complete-consistent-sets

\documentclass[../../../include/open-logic-section]{subfiles}

\begin{document}

\olfileid[id]{nml}{com}{ccs}

\olsection{Himpunan Lengkap dan $\Sigma$-Konsisten}

Andaikan $\Sigma$ merupakan suatu himpunan !!{formula}s modal---anggaplah
formula-formula itu sebagai aksioma atau prinsip pendefinisi suatu logika
modal normal. Suatu himpunan $\Gamma$ $\Sigma$-konsisten jika dan hanya jika
$\Gamma \Proves/[\Sigma] \lfalse$, yakni jika tidak terdapat !!{derivation}
dari~$!A_1 \lif (!A_2 \lif \cdots (!A_n \lif \lfalse)\dots)$ berdasarkan
$\Sigma$, dengan setiap $!A_i \in \Gamma$. Kita akan membangun suatu model
``kanonik'' yang setiap dunianya diambil sebagai jenis khusus himpunan
$\Sigma$-konsisten: himpunan yang bukan hanya $\Sigma$-konsisten, melainkan
maksimal demikian, dalam arti bahwa himpunan tersebut menentukan nilai
kebenaran setiap !!{formula} modal: untuk setiap $!A$, berlaku $!A \in
\Gamma$ atau $\lnot !A \in \Gamma$:

\begin{defn}
  Suatu himpunan $\Gamma$ \emph{lengkap dan $\Sigma$-konsisten} jika dan
  hanya jika himpunan itu $\Sigma$-konsisten dan, untuk setiap~$!A$, berlaku
  $!A \in \Gamma$ atau $\lnot !A \in \Gamma$.
\end{defn}

Himpunan lengkap dan $\Sigma$-konsisten~$\Gamma$ mempunyai sejumlah sifat
yang berguna. Salah satunya, himpunan tersebut tertutup secara deduktif,
yakni jika $\Gamma \Proves[\Sigma] !A$, maka $!A \in \Gamma$. Secara khusus,
ini berarti bahwa setiap instansiasi dari !!a{formula}~$!A \in \Sigma$ juga
memenuhi $!A \in \Gamma$. Selain itu, keanggotaan dalam $\Gamma$
mencerminkan syarat kebenaran bagi penghubung proposisional. Hal ini akan
penting ketika kita mendefinisikan ``model kanonik.''

\begin{prop}\ollabel{prop:ccs-properties}
  Andaikan $\Gamma$ lengkap dan $\Sigma$-konsisten. Maka:
  \begin{enumerate}
  \item \ollabel{prop:ccs-closed}%
    $\Gamma$ tertutup secara deduktif dalam~$\Sigma$.
  \item \ollabel{prop:ccs-sigma}%
    $\Sigma \subseteq \Gamma$.
  \tagitem{prvFalse}{\ollabel{prop:ccs-lfalse}%
    $\lfalse \notin \Gamma$}{}
  \tagitem{prvTrue}{\ollabel{prop:ccs-ltrue}%
    $\ltrue \in \Gamma$}{}
  \item \ollabel{prop:ccs-lnot}%
    $\lnot!A \in \Gamma$ jika dan hanya jika $!A \notin
    \Gamma$.
  \tagitem{prvAnd}{\ollabel{prop:ccs-land}%
    $!A \land !B \in \Gamma$ jika dan hanya jika $!A \in \Gamma$ dan $!B \in \Gamma$}{}
  \tagitem{prvOr}{\ollabel{prop:ccs-lor}%
    $!A \lor !B \in \Gamma$ jika dan hanya jika $!A \in \Gamma$ atau $!B \in \Gamma$}{}
  \tagitem{prvIf}{\ollabel{prop:ccs-lif}%
    $!A \lif !B \in \Gamma$ jika dan hanya jika $!A \notin \Gamma$ atau $!B \in \Gamma$}{}
  \tagitem{prvIff}{\ollabel{prop:ccs-liff}%
    $!A \liff !B \in \Gamma$ jika dan hanya jika $!A \in \Gamma$ dan $!B \in
    \Gamma$, atau $!A \notin \Gamma$ dan $!B \notin \Gamma$}{}
  \end{enumerate}
\end{prop}

\begin{proof}
  \begin{enumerate}
  \item Andaikan $\Gamma \Proves[\Sigma] !A$, tetapi $!A \notin
    \Gamma$. Karena $\Gamma$ lengkap dan $\Sigma$-konsisten,
    $\lnot!A \in \Gamma$. Hal ini akan membuat $\Gamma$ tak konsisten,
    sebab $!A, \lnot !A \Proves[\Sigma] \lfalse$.

  \item Jika $!A \in \Sigma$, maka $\Gamma \Proves[\Sigma] !A$, dan $!A
    \in \Gamma$ berdasarkan ketertutupan deduktif, yakni
    kasus~\olref{prop:ccs-closed}.

  \tagitem{prvFalse}{Jika $\lfalse \in \Gamma$, maka $\Gamma
    \Proves[\Sigma] \lfalse$, sehingga $\Gamma$ akan menjadi
    tidak $\Sigma$-konsisten.}{}

  \tagitem{prvTrue}{$\Gamma \Proves[\Sigma] \ltrue$, sehingga $\ltrue \in
    \Gamma$ berdasarkan ketertutupan deduktif, yakni
    kasus~\olref{prop:ccs-closed}.}{}

\item Jika $\lnot!A \in \Gamma$, maka berdasarkan konsistensi $!A \notin
    \Gamma$; dan jika $!A \notin \Gamma$, maka $\lnot!A \in \Gamma$ karena
    $\Gamma$ lengkap dan $\Sigma$-konsisten.

  \tagitem{prvAnd}{\iftag{probAnd}{Latihan.}{Andaikan $!A \land !B \in
      \Gamma$. Karena $(!A \land !B) \lif !A$ merupakan instansiasi
      tautologis, $!A \in \Gamma$ berdasarkan ketertutupan deduktif, yakni
      kasus~\olref{prop:ccs-closed}. Demikian pula untuk $!B \in \Gamma$.
      Sebaliknya, andaikan $!A \in \Gamma$ dan $!B \in \Gamma$. Maka
      ketertutupan deduktif mengakibatkan $(!A \land !B) \in \Gamma$, sebab
      $!A \lif (!B \lif (!A \land !B))$ merupakan instansiasi tautologis.}}{}

  \tagitem{prvOr}{\iftag{probOr}{Latihan.}{Andaikan $!A \lor !B \in
      \Gamma$, serta $!A \notin \Gamma$ dan $!B \notin \Gamma$. Karena
      $\Gamma$ lengkap dan $\Sigma$-konsisten, $\lnot!A \in \Gamma$ dan
      $\lnot !B \in \Gamma$. Maka $\lnot(!A \lor !B) \in \Gamma$ sebab
      $\lnot !A \lif (\lnot !B \lif \lnot (!A \lor !B))$ merupakan
      instansiasi tautologis. Hal ini berarti bahwa $\Gamma$ tak
      $\Sigma$-konsisten, suatu kontradiksi.}}{}

  \tagitem{prvIf}{\iftag{probIf}{Latihan.}{Andaikan $!A \lif !B \in
      \Gamma$ dan $!A \in \Gamma$; maka $\Gamma \Proves[\Sigma] !B$,
      sehingga $!B \in \Gamma$ berdasarkan ketertutupan deduktif.
      Sebaliknya, jika $!A \lif !B \notin \Gamma$, maka karena $\Gamma$
      lengkap dan $\Sigma$-konsisten, $\lnot (!A \lif !B) \in \Gamma$.
      Karena $\lnot(!A \lif !B) \lif !A$ merupakan instansiasi tautologis,
      $!A \in \Gamma$ berdasarkan ketertutupan deduktif. Karena
      $\lnot(!A \lif !B) \lif \lnot !B$ merupakan instansiasi tautologis,
      $\lnot !B \in \Gamma$. Maka $!B \notin \Gamma$ karena $\Gamma$
      $\Sigma$-konsisten.}}{}

  \tagitem{prvIff}{\iftag{probIff}{Latihan.}{Andaikan $!A \liff !B \in
      \Gamma$. Jika $!A \in \Gamma$, maka $!B \in \Gamma$, sebab $(!A
      \liff !B) \lif (!A \lif !B)$ merupakan instansiasi tautologis.
      Demikian pula, jika $!B \in \Gamma$, maka $!A \in \Gamma$. Jadi,
      berlaku sekaligus $!A \in \Gamma$ dan $!B \in \Gamma$, atau tidak
      berlaku baik $!A \in \Gamma$ maupun $!B \in \Gamma$.

      Sebaliknya, andaikan $!A \liff !B \notin \Gamma$. Karena $\Gamma$
      lengkap dan $\Sigma$-konsisten, $\lnot (!A \liff !B) \in \Gamma$.
      Karena $\lnot(!A \liff !B) \lif (!A \lif \lnot !B)$ merupakan
      instansiasi tautologis, jika $!A \in \Gamma$, maka $\lnot !B \in
      \Gamma$, dan karena $\Gamma$ $\Sigma$-konsisten, $!B \notin \Gamma$.
      Demikian pula, jika $!B \in \Gamma$, maka $!A \notin \Gamma$.
      Jadi, tidak berlaku bahwa sekaligus $!A \in \Gamma$ dan $!B \in
      \Gamma$, dan juga tidak berlaku bahwa sekaligus $!A \notin \Gamma$
      dan $!B \notin \Gamma$.}}{}
  \end{enumerate}
\end{proof}

\begin{probtag}{probAnd,probOr,probIf,probIff}
  Lengkapi pembuktian \olref[nml][com][ccs]{prop:ccs-properties}.
\end{probtag}

\end{document}
